\documentclass[11pt]{article}
\usepackage[margin=1in]{geometry}
\usepackage{amsmath,amssymb}
\usepackage{booktabs}
\usepackage{hyperref}

\title{Hierarchical Markov Modeling for Underwater Pursuit-Evasion}
\author{}
\date{}

\begin{document}
\maketitle

\section{Overview}
This document provides three consistent formulations based on the project implementation:
\begin{itemize}
  \item Low-level 6DOF control MDP.
  \item High-level pursuit-evasion task MDP.
  \item Macro-action (Options) SMDP for hierarchical control.
\end{itemize}

\section{Low-Level 6DOF Control MDP}
We model the low-level controller as an MDP:
\begin{equation}
M_{LL} = (\mathcal{S}_{LL}, \mathcal{A}_{LL}, P_{LL}, R_{LL}, \gamma).
\end{equation}

\subsection*{State}
The state is the full 6DOF vehicle state plus goal:
\begin{equation}
 s_t^{LL} = (\eta_t, \nu_t, g_t),
\end{equation}
where $\eta_t = [x, y, z, \phi, \theta, \psi]^\top$ is position and orientation, $\nu_t = [u, v, w, p, q, r]^\top$ is linear and angular velocity, and $g_t$ denotes the target state (position/orientation, and optionally target velocity).

\subsection*{Action}
The low-level action space is discrete with 13 semantic actions:
\begin{equation}
\mathcal{A}_{LL} = \{a_0, a_1, \ldots, a_{12}\}.
\end{equation}
They correspond to hover, surge/sway/heave, and yaw/pitch/roll in positive or negative directions.

\subsection*{Transition}
The action is mapped to thrust through a control stack, then the 6DOF dynamics evolve:
\begin{align}
\tau_t &= f_{ctrl}(a_t, s_t^{LL}), \\
 s_{t+1}^{LL} &= f_{dyn}(s_t^{LL}, \tau_t, \xi_t),
\end{align}
where $\xi_t$ denotes disturbances (water current, domain randomization).

\subsection*{Reward}
The reward combines position/orientation error, velocity penalty, energy cost, and goal/time terms:
\begin{equation}
 r_t^{LL} = w_p\|e_p\| + w_o\|e_o\| + w_v\|\nu\| + w_e\|\tau\| + r_{goal} + r_{time}.
\end{equation}

\subsection*{Termination}
Episodes terminate on success (within thresholds), out-of-bounds, flip, no-progress, or timeout.

\section{High-Level Pursuit-Evasion MDP}
The pursuit-evasion task is modeled as:
\begin{equation}
M_H = (\mathcal{S}_H, \mathcal{A}_H, P_H, R_H, \gamma_H).
\end{equation}

\subsection*{State}
A task-level state can be defined by relative geometry and boundary information:
\begin{equation}
 s_t^H = [p_p, v_p, p_e, v_e, \psi_e, b_p, b_e],
\end{equation}
where $p_p, v_p$ are pursuer position and velocity, $p_e, v_e$ are evader position and velocity, $\psi_e$ is evader heading, and $b_p, b_e$ represent boundary margins.

\subsection*{Action}
High-level actions represent sub-goal selection or pursuit strategies:
\begin{equation}
 \mathcal{A}_H = \{\text{subgoal/heading/strategy}\}.
\end{equation}

\subsection*{Transition}
High-level actions are realized by running the low-level policy for multiple steps:
\begin{equation}
 s_{t+1}^{H} = F(s_t^{H}, a_t^{H}, \pi_{LL}).
\end{equation}

\subsection*{Reward}
A typical high-level reward for pursuit-evasion:
\begin{equation}
 r_t^H = \alpha\Delta d + \beta v_{close} + r_{catch} + r_{boundary} - r_{time}.
\end{equation}

\subsection*{Termination}
Capture, failure (escape/out-of-bounds), or timeout.

\section{Macro-Action / Options SMDP}
The hierarchical control can be modeled as an SMDP with options:
\begin{equation}
M_S = (\mathcal{S}, \mathcal{O}, P_S, R_S, \gamma).
\end{equation}

\subsection*{Options}
Each option is defined as:
\begin{equation}
 o = (\mathcal{I}_o, \pi_o, \beta_o),
\end{equation}
where $\mathcal{I}_o$ is the initiation set, $\pi_o$ is the intra-option policy (low-level controller), and $\beta_o$ is the termination condition.

\subsection*{Transition}
An option runs for $k$ steps before transitioning:
\begin{equation}
 P_S(s'\mid s,o) = \sum_{k\ge 1} P(s', k \mid s, o).
\end{equation}

\subsection*{Return}
The SMDP reward accumulates over the option duration:
\begin{equation}
 R_S(s,o) = \mathbb{E}\left[\sum_{i=0}^{k-1}\gamma^i r_{t+i}\right].
\end{equation}

\section{Observation Vector Definitions}
This section enumerates the exact observation vectors used in the code.

\subsection{2D Pursuit-Evasion Observation (22D)}
Let $\Delta p = p_e - p_p$, $\Delta v = v_e - v_p$, $d = \|\Delta p\|$, and
\begin{equation}
 v_{close} = -\frac{\Delta v \cdot \Delta p}{\|\Delta p\|} \quad (d>0).
\end{equation}
Normalization constants:
\begin{equation}
\text{pos\_scale}=\frac{1}{\text{world\_size}},\quad
\text{vel\_scale}=\frac{1}{\max(v_{p\_max}, v_{e\_max})},\quad
\text{acc\_scale}=\frac{1}{\max(a_{p\_max}, a_{e\_max})}.
\end{equation}
The 22-dimensional observation is ordered as:
\begin{align*}
&1: p_{p,x}\cdot\text{pos\_scale}, \quad
2: p_{p,y}\cdot\text{pos\_scale}, \\
&3: v_{p,x}\cdot\text{vel\_scale}, \quad
4: v_{p,y}\cdot\text{vel\_scale}, \\
&5: \Delta p_x\cdot\text{pos\_scale}, \quad
6: \Delta p_y\cdot\text{pos\_scale}, \\
&7: \Delta v_x\cdot\text{vel\_scale}, \quad
8: \Delta v_y\cdot\text{vel\_scale}, \\
&9: d\cdot\text{pos\_scale}, \quad
10: v_{close}\cdot\text{vel\_scale}, \\
&11: \hat{v}_{e,x}, \quad
12: \hat{v}_{e,y}, \\
&13: (\text{world\_size} - p_{p,y})\cdot\text{pos\_scale}, \\
&14: (p_{p,y} + \text{world\_size})\cdot\text{pos\_scale}, \\
&15: (p_{p,x} + \text{world\_size})\cdot\text{pos\_scale}, \\
&16: (\text{world\_size} - p_{p,x})\cdot\text{pos\_scale}, \\
&17: \min(\text{world\_size}-|p_{e,x}|,\;\text{world\_size}-|p_{e,y}|)\cdot\text{pos\_scale}, \\
&18: \|v_p\|\cdot\text{vel\_scale}, \quad
19: \|v_e\|\cdot\text{vel\_scale}, \\
&20: \hat{v}_p\cdot\hat{\Delta p}, \\
&21: a_{p,x}\cdot\text{acc\_scale}, \quad
22: a_{p,y}\cdot\text{acc\_scale}.
\end{align*}

\subsection{Low-Level 6DOF Observation}
Let $R(\phi,\theta,\psi)$ be the rotation matrix from body to world. The body-frame position error is:
\begin{equation}
 e_p^{body} = R(\phi,\theta,\psi)^\top (g_p - \eta_{pos}).
\end{equation}
The orientation error is the wrapped angle difference:
\begin{equation}
 e_o = \text{wrap}(\eta_{ori} - g_{ori}) \in [-\pi,\pi]^3.
\end{equation}
The observation is assembled in segments depending on configuration flags. Each segment is appended in the following order:
\begin{itemize}
  \item Position error (body frame): $e_p^{body}/10$.
  \item Orientation error: $e_o/\pi$.
  \item Linear velocity (body frame): $[u,v,w]/2$.
  \item Angular velocity (body frame): $[p,q,r]/2$.
  \item Target position (body frame): $e_p^{body}/10$.
  \item Target orientation: $g_{ori}/\pi$.
  \item Error terms: $e_p^{body}/10$ and $e_o/\pi$.
\end{itemize}
The final observation is the concatenation of all enabled segments in that order.

\section{How to Compile}
From the project root, run:
\begin{verbatim}
cd docs
pdflatex markov_modeling.tex
\end{verbatim}

\end{document}
